\begin{abstract}
Visual tool-use distillation should teach models how to acquire evidence, not merely how to produce correct final answers. We present \method, a multi-domain trajectory dataset and synthesis recipe for visual tool use. \dataset contains 42K trajectories across Chart, GUI Grounding, Table, Visual Search, and Web-to-HTML tasks, covering complementary visual bottlenecks such as small text, dense layouts, coordinate precision, object search, color or structure analysis, and render-based verification. The key design principle is that a retained trajectory must satisfy two conditions: the task should benefit from tools, and the recorded tool outputs should actually improve solving the task. To enforce this, our pipeline first removes examples that are easily solved without tools, then distills teacher tool-use trajectories, checks final-answer correctness, and filters for \toolgain by measuring whether the teacher's tool observations improve a base model's performance. Fine-tuning on \dataset improves all five domains across four backbones, with average gains of 8.9--13.2 points. Ablations show that correctness and \toolgain filtering are complementary, and comparison with no-tool trajectory SFT shows that tool-use supervision is especially beneficial for evidence-acquisition tasks such as Visual Search and Table.
\end{abstract}

\section{Introduction}
\label{sec:introduction}

Many visual tasks are hard not because the reasoning is abstract, but because the relevant evidence is hard to observe in the original input: reading small axis labels in a chart, aligning a row and column across a low-contrast table, locating a tiny control in a high-resolution screenshot, checking small or occluded objects in a cluttered scene, or rendering generated HTML to verify that its layout matches a reference. In all of these cases the model should actively inspect, transform, mark, or verify the visual input rather than answer from a single static image---the paradigm of \emph{thinking with images} \citep{openai2025o3,su2025thinkingsurvey}, in which a model acts on the visual input as part of its reasoning chain rather than reasoning over a fixed view.

Teaching this behavior through supervised distillation requires trajectories that record the evidence-acquisition process: what the model chose to inspect, what the tool returned, and how the new observation changed its subsequent reasoning. This is exactly where building such data is hard. A dataset that records only final answers misses the process, while one that retains arbitrary tool calls teaches the model to imitate tool-use syntax rather than to use tools when they help. The standard recipe---distill from a strong teacher and keep a trajectory whenever its final answer is correct---does not resolve this tension: a teacher can reach the right answer through unnecessary tool calls, and final-answer correctness reveals neither whether the task needed tools nor whether the recorded observations actually helped.

We therefore argue that a retained trajectory must satisfy two conditions: the task should benefit from tools, and the recorded tool outputs should measurably improve solving it. We operationalize both with \method, an open recipe that curates trajectories by double rejection sampling. We first discard queries a base model already solves without tools, so that external evidence has room to help. We then roll out a strong teacher under a unified visual toolset and keep a trajectory only when it is both \emph{answer-correct}, verified by domain-appropriate evaluators (rule-based matching, an LLM judge, or render-based VLM judging for Web-to-HTML), and \emph{tool-gainful}: prepending the teacher's tool observations to a base model's context must improve its accuracy on the query. The main \dataset is the intersection of the two sets.

The resulting \dataset contains 42K trajectories spanning five domains chosen for complementary visual bottlenecks: Chart and Table stress fine-grained reading and structural alignment, GUI Grounding stresses local inspection and coordinate precision, Visual Search stresses active evidence search in cluttered scenes, and Web-to-HTML stresses a generate-render-observe-revise loop. All domains share one visual toolset---cropping, marking, color and structure analysis, rendering, and final action expression---so the model learns reusable operations rather than a separate interface per task.

Fine-tuning on \dataset improves all five domains across four backbones spanning two model families and 4B--27B scale. On Qwen3.5-9B the average rises from 32.6 to 45.8, with the largest gains where tools expose otherwise-hidden evidence (Visual Search +23.8, Table +17.3). Controlled variants confirm the design: correctness-only and tool-gain-only filtering each underperform their intersection on every domain, and against a no-tool trajectory SFT baseline trained on the same queries, tool-use supervision wins on four of five domains and by 8.7 average points. The exception, GUI Grounding, is instructive rather than hidden: when the metric rewards a single coordinate, direct imitation can be more sample-efficient than a multi-step protocol of cropping, coordinate remapping, and action formatting. Visual tool-use distillation is thus not a uniform replacement for all supervision; its advantage is sharpest when tools acquire new visual evidence.

In summary, our contributions are:
\begin{itemize}
  \item We present \method, an open pipeline for synthesizing multi-domain visual tool-use trajectories under a shared visual toolset.
  \item We construct \dataset, a 42K-example corpus covering Chart, GUI Grounding, Table, Visual Search, and Web-to-HTML, chosen for complementary visual bottlenecks and tool demands.
  \item We propose double rejection sampling based on task difficulty, final-answer correctness, and \toolgain, retaining only trajectories that both solve the task and contain tool observations that improve solvability.
  \item We provide controlled dataset variants and experiments showing that the correctness and \toolgain filters are complementary, that tool-use trajectories outperform no-tool trajectories on evidence-acquisition domains, and that GUI Grounding exposes a useful limitation of the current action protocol.
\end{itemize}
